\section{Evolução do Esquema}
\label{sec:evolucao}

O esquema muda ao longo do projeto: cada funcionalidade nova pode pedir uma
tabela ou uma coluna. Essas mudanças são feitas por \textbf{migrations} do
Flyway, arquivos SQL aplicados em ordem e registrados na tabela
\campo{flyway\_schema\_history}. O Flyway roda a cada \campo{docker compose
up}, então o banco de qualquer integrante e o de produção ficam sempre na
mesma versão.

\subsection{Nome e ordem}

O nome de cada arquivo segue o padrão
\campo{V<AAMMDD>\_<HHMMSS>\_\_<Descricao>.sql}. Como a versão é a data e a hora
de criação, as migrations de integrantes diferentes não disputam o mesmo
número, e rodam na ordem em que foram escritas. Um nome fora do padrão faz o
Flyway falhar (\campo{validateMigrationNaming}), em vez de ignorar o arquivo em
silêncio.

\begin{table}[htbp]
  \centering
  \caption{Migrations aplicadas até a entrega}
  \label{tab:migrations}
  \small
  \begin{tabularx}{\textwidth}{@{}L{0.42}L{1.58}@{}}
    \toprule
    \poliCabTabela{Versão} & \poliCabTabela{O que muda} \\
    \midrule
    \campo{V1} & Esquema inicial: planos, usuários, catálogo, interações e playlists \\
    \campo{V261020\_162546} & Tabelas do Better Auth e email na tabela \campo{user} \\
    \campo{V261026\_133836} & Papel do usuário (\campo{admin} ou \campo{assinante}) \\
    \campo{V261103\_151718} & Extensão \campo{unaccent}, para a busca não diferenciar acentos \\
    \campo{V261103\_154132} & Eventos de reprodução passam a apontar para a música \\
    \campo{V261103\_163300} & Ordem das músicas na playlist \\
    \campo{V261105\_151000} & Convites do plano familiar e limite de 4 pessoas \\
    \campo{V261105\_162500} & Denúncias de comentários \\
    \campo{V261105\_183723} & Áudio e conteúdo explícito na música, ano no álbum \\
    \campo{V261105\_213000} & Capa do álbum \\
    \campo{V261112\_120000} & Número da faixa no álbum \\
    \campo{V261115\_200000} & Playlist ``Músicas Curtidas'' de cada conta \\
    \campo{V261116\_090000} & Evento \campo{music\_listened} \\
    \campo{V261116\_090100} & O evento \campo{music\_listened} também leva a música \\
    \campo{V261117\_160000} & Nome de artista único, com a fusão dos duplicados \\
    \campo{V261117\_170000} & Capa da playlist e foto do artista \\
    \bottomrule
  \end{tabularx}
\end{table}

\FloatBarrier
\subsection{Regras}

\begin{itemize}
  \item \textbf{Uma migration publicada nunca é editada.} O Flyway guarda um
    checksum de cada arquivo aplicado e se recusa a subir se ele mudar. Uma
    correção é sempre uma migration nova.
  \item \textbf{Dados de teste não entram em migration}, porque as migrations
    também rodam em produção. Eles ficam no seed (\cref{sec:acesso}).
  \item \textbf{Uma regra nova cuida dos dados que já existem.} Antes de criar
    o índice de nome de artista único, a migration
    \campo{V261117\_160000} funde os artistas de mesmo nome no de menor
    \campo{id}: os vínculos com álbuns e músicas passam para ele, e os outros
    recebem remoção lógica. Sem isso, o índice falharia em bancos que já
    tivessem duplicados.
  \item \textbf{Valor novo de enum e o uso dele vão em migrations separadas.}
    O Postgres não permite usar um valor de enum na mesma transação em que ele
    foi criado, e o Flyway roda cada migration numa transação. Por isso
    \campo{music\_listened} entrou numa migration, e o \campo{CHECK} que o usa
    entrou na seguinte.
\end{itemize}

\subsection{Documentação gerada}

Depois de uma migration, o comando \campo{db:sync} atualiza o
\campo{schema.prisma} a partir do banco, e o \campo{db:sync --check} confere se
os dois estão iguais. O diagrama da \cref{fig:relacional} é gerado pelo CI.
Com isso, a documentação do esquema não depende de alguém lembrar de
atualizá-la.
